\section{Group-Specific Fixed Execution}
\label{sec:execution}

The executor maintains one action-sequence buffer $B^{(g)}$, cursor $c_g$, and
remaining counter $r_g$ for each group.  If any $r_g=0$, the current observation
is passed once through the frozen policy to obtain a fresh full chunk.  Every
expired group copies its slice of that chunk into $B^{(g)}$, resets $c_g=0$, and
sets $r_g=h_g$.  A non-expired group keeps both its buffer and cursor.  The
executed action is then assembled as
\begin{equation}
\mathbf a_t=
\bigl[B^{(1)}[c_1],\ldots,B^{(G)}[c_G]\bigr],
\label{eq:compose}
\end{equation}
after which all cursors advance and all remaining counters decrease.

This rule has three consequences.  First, a policy query always predicts a full
joint chunk, even if only one group accepts it; query rate therefore measures
full policy evaluations, not accepted components.  Second, non-expired groups
continue along their own previously predicted sequence rather than holding their
last action.  Third, components in (\ref{eq:compose}) can have different source
observations and group-local sequence positions.

Starting with all groups expired at step 0, the nominal query-boundary set for a
fixed-length episode is
\begin{equation}
\mathcal Q(\mathbf h)=
\bigcup_{g=1}^{G}\{k h_g\mid k=0,1,2,\ldots\}.
\label{eq:query_set}
\end{equation}
If several groups expire together, their update uses one shared query.  Swapping
the same horizon values between groups leaves $\mathcal Q$ unchanged.  The
symmetric comparisons in Sec.~\ref{sec:results} therefore have the same nominal
query schedule; small differences in aggregate query rate arise only because
configurations can terminate episodes at different steps.  For the evaluated
powers-of-two pairs, the shorter horizon divides the longer and the query
schedule is set by the shorter one.  This property makes horizon swaps a direct
test of assignment rather than policy-call frequency.

Fig.~\ref{fig:timeline} shows the exact schedule for global $h=4$ and
$(h_a,h_g)=(4,16)$.  At steps 4, 8, and 12, the policy produces a fresh full
chunk, but only the expired arm group accepts it; the gripper continues chunk
$C0$.  At step 16 both groups expire and accept $C4$.  This mechanism tests
independent horizon assignment without modifying policy weights, but it does not
guarantee cross-group temporal consistency.

\begin{figure*}[t]
    \centering
    \includegraphics[width=0.96\textwidth]{execution_timeline.pdf}
    \caption{Fixed execution schedules.  Colors and labels denote the policy-query
    generation supplying each group.  Global execution replaces all components
    together.  With $(h_a,h_g)=(4,16)$, arm refreshes trigger full policy queries,
    but the non-expired gripper retains its earlier sequence until step 16.}
    \label{fig:timeline}
\end{figure*}
